\section{Overall Description}

\subsection{Product Perspective}
The JUST-DAPM-PGC is a web-based automated admission management system designed to replace the existing manual and partially digital processes. The previous system lacked critical features such as proper validation mechanisms, deadline enforcement, and structured system architecture.

\medskip
\noindent The proposed system introduces:
\begin{itemize}[leftmargin=1.5cm, labelsep=0.5cm]
    \item Strong input validation to prevent invalid or incomplete submissions.
    \item Deadline control to restrict registration and updates after the admission deadline.
    \item Role-based access control to ensure secure and authorized operations.
    \item A modular and maintainable architecture following modern design patterns.
    \item Real-time tracking of application progress across different offices.
    \item Scalability to handle a large number of concurrent applicants.
\end{itemize}

\subsection{Product Functions}
The major functionalities of the JUST-DAPM-PGC are as follows:

\begin{itemize}[leftmargin=0.5cm]
    \item[] Admin Module: 
    This module takes care of administration activities such as assigning different roles to users and setting the system security features. Administrators can keep track of the status of the admissions process through the use of this module.
    \begin{itemize}
        \item Manage users and roles.
        \item Control system operations (open/close admission).
        \item Monitor overall admission progress.
    \end{itemize}
    
    \vspace{0.2cm}
    \item[] Dean Office Module: 
    Charged with the responsibility of evaluating the academic qualifications of students through verification of education certificates. The Dean’s office makes it possible to either approve or disapprove applications depending on individual faculty requirements.
    \begin{itemize}
        \item Verify documents and validate applicants.
        \item Approve or reject applications.
        \item Export applicant data.
    \end{itemize}
    
    \vspace{0.2cm}
    \item[] Hall Office Module: 
    Deals with all financial and residence related issues such as the clearance of payment of admission costs using bank slips. This module ensures that all required residence clearance activities are completed for students requiring on-campus accommodation.
    \begin{itemize}
        \item Verify bank receipts.
        \item Manage residential clearance.
    \end{itemize}
    
    \vspace{0.2cm}
    \item[] Medical Center Module: 
    Enables the health evaluation process by capturing applicant’s medical details and carrying out physical examinations. This module creates medical reports which provide required health clearance before completing the admissions process.
    \begin{itemize}
        \item Record medical data.
        \item Generate reports and provide clearance.
    \end{itemize}
    
    \vspace{0.2cm}
    \item[] Registrar Office Module: 
    The last resort in the approval of applications for admissions. The module collects, stores and manages the information of admitted students through archiving of documents.
    \begin{itemize}
        \item Final admission approval.
        \item Document collection and archiving.
    \end{itemize}
    
    \vspace{0.2cm}
    \item[] Applicant Module: 
    The module helps the students register themselves, update their profile information, and apply for admissions by submitting the application form. The applicant can upload all relevant documents and even download important cards/notifications.
    \begin{itemize}
        \item Registration and profile management.
        \item Application form submission.
        \item Document upload and download.
    \end{itemize}
\end{itemize}